\documentclass[12pt]{article}
\usepackage[margin=1in]{geometry}
\usepackage{amsmath}

\begin{document}

\section*{Figure Captions}

\noindent\textbf{Figure 1. Multiple Equilibria.}

\noindent\textit{Notes:} The blue curve shows the utility advantage of green work relative to non-green work as the share of green workers increases. The horizontal zero line marks indifference; crossings and dots are equilibria.

\bigskip

\noindent\textbf{Figure 2. Two policy levers to raise green employment.}

\noindent\textit{Panel (a):} Lower worker entry cost raises the green share.  

\noindent\textit{Panel (b):} Higher firm subsidy raises the green share.  

\noindent\textit{Notes:} Axes are the same as in Figure 1. Comparative statics are with respect to each policy margin, holding other parameters fixed.

\bigskip

\noindent\textbf{Figure 3. Green Talent Demand and Supply (Indexed; Projections to 2050).}

\noindent\textit{Notes:} The figure plots global indexes of green talent demand (dark line) and green talent supply (light line). Both series are normalized to 1 in 2021; projections from 2025 onward use median values and compound annual growth rates. By 2050, demand is about 5.0 and supply about 2.5, indicating a widening gap. The y-axis reports index values (2021 = 1). \textit{Source:} LinkedIn (2024) \textit{Green Skills and Jobs} stocktake report.

\bigskip

\noindent\textbf{Figure 4. Worker flows between unemployment and employment in both sectors.}

\noindent\textit{Notes:} The diagram summarizes flows between unemployment and employment in the green and non-green sectors. Unemployed workers choose a sector to search in; matches move workers into employment, and separations return workers to unemployment at rate $\lambda$.

\bigskip

\noindent\textbf{Figure 5. Dynamic model: two policy instruments to increase green employment.}

\noindent\textit{Panel (a):} Lower worker entry costs raise the green jobs share.  

\noindent\textit{Panel (b):} Higher firm subsidies raise the green jobs share.  

\noindent\textit{Notes:} Panels show dynamic comparative statics holding other parameters fixed. Lower worker entry costs or higher firm subsidies increase the share of green jobs.

\bigskip

\noindent\textbf{Figure 6. Welfare gains as a function of the externality parameter $\phi$ ($\phi^\star \approx 0.062$).}

\noindent\textit{Notes:} The blue line shows welfare with the green-job externality minus baseline welfare. The horizontal zero line marks equal welfare. The dashed vertical line indicates the threshold $\phi^\star \approx 0.062$; welfare is higher than baseline for $\phi > \phi^\star$ and lower for $\phi < \phi^\star$.

\bigskip

\noindent\textbf{Figure 7. Funding decomposition along the transition path.}

\noindent\textit{Notes:} Bars split total funding into worker-side and firm-side components; the line overlays total funding relative to GDP. All figures are computed along the welfare-maximizing transition to a 14\% green employment share.

\bigskip

\noindent\textbf{Figure 8. Green labor market dynamics along the transition path.}

\noindent\textit{Panel (a):} Green market tightness, $\theta_g \equiv v_g/u_g$, along the welfare-maximizing path. Tightness initially rises as vacancies outpace entrants, peaks when the market becomes very tight, and then declines as policy tilts back toward worker entry.  

\noindent\textit{Panel (b):} Job-finding $f(\theta_g) = \text{match}_g/u_g$ and vacancy-filling $q(\theta_g) = \text{match}_g/v_g$ along the same path. As $\theta_g$ rises, $f(\theta_g)$ increases and then flattens mechanically, while $q(\theta_g)$ declines and reaches a trough; when $\theta_g$ later falls, $q(\theta_g)$ recovers.

\bigskip

\noindent\textbf{Figure 9. Pathwise sensitivities of matching rates to tightness.}

\noindent\textit{Notes:} Sensitivities are computed as finite differences along the welfare-maximizing sequence of steady states: $\Delta f/\Delta\theta_g$ and $\Delta q/\Delta\theta_g$. Early in the path, the job-finding rate $f(\theta_g)$ is most responsive to changes in tightness while vacancy-filling $q(\theta_g)$ remains relatively high; mid-path, both margins are locally inelastic; late in the path, $q(\theta_g)$ becomes more responsive as $\theta_g$ recedes.

\bigskip
\bigskip

\section*{Table Captions}

\noindent\textbf{Table 1. Model performance in matching calibration targets.}

\noindent\textit{Notes:} The table compares model-implied moments to empirical targets used in the calibration, including the green wage premium, green employment share, hiring likelihood ratio, aggregate labor market tightness, and the unemployment rate.

\bigskip

\noindent\textbf{Table 2. Internally calibrated parameters.}

\noindent\textit{Notes:} The table reports parameter values chosen to match calibration targets, including the vacancy creation cost $c$, separation rate $\lambda$, sector-specific matching efficiencies $(\delta_g,\delta_n)$, and the green worker entry cost $\kappa_g$.

\bigskip

\noindent\textbf{Table 3. Alternative policy mixes achieving 14\% green employment.}

\noindent\textit{Notes:} The table reports alternative policy configurations (dual intervention, worker-only, firm-only) that all reach an equilibrium green employment share of 14\%, starting from a 2\% baseline. For each configuration, it shows worker-side subsidies, firm-side subsidies, total subsidies (as a percentage of GDP), and the resulting green employment share.

\bigskip

\noindent\textbf{Table 4. Policy trade-offs: unemployment, fiscal cost, and welfare.}

\noindent\textit{Notes:} The table compares unemployment, fiscal cost (as a percentage of GDP), and welfare (consumption-equivalent variation relative to baseline) across the dual, worker-only, and firm-only policy configurations, conditional on all reaching 14\% green employment. The welfare measure incorporates a calibrated green-job externality $\phi$ that yields a 2\% consumption-equivalent gain in the benchmark case; the baseline (2022) welfare level is normalized to zero.

\end{document}
